\chapter{Introduction}

The Linear Ordering Problem (LOP) is a well-known NP-hard combinatorial
optimisation problem with a wide range of practical applications, including the
triangularisation of input-output matrices in economics, chronological ordering
in archaeology, and scheduling problems~\cite{SchStu04, MarReiDua12}.
Given the difficulty of finding optimal solutions in reasonable time, heuristic
and metaheuristic approaches are commonly used in practice.

Local search methods form the backbone of many successful algorithms for the
LOP. They operate by iteratively moving from a current solution to a
neighbouring one as long as an improvement can be found, stopping at a local
optimum. The quality of the final solution largely depends on the choice of
neighbourhood structure and pivoting rule used during the search.

This report describes the implementation and experimental evaluation of
iterative improvement (II) algorithms and Variable Neighbourhood Descent (VND)
for the LOP. For II, we consider three neighbourhood structures (transpose,
exchange, and insert) combined with two pivoting rules (first-improvement
and best-improvement) and two initialisation strategies (random permutation
and the Chenery-Watanabe greedy heuristic~\cite{CheWat58}), yielding twelve
algorithm configurations in total. VND extends this by cycling through
multiple neighbourhoods to escape local optima that a single neighbourhood
cannot avoid~\cite{HanMla01}.

A key aspect of the implementation is the use of incremental delta evaluations
for all neighbourhood structures, which avoids recomputing the full objective
function at each step and significantly reduces computation time.

The algorithms are tested on 78 benchmark instances of sizes 150 and 250,
drawn from established LOP benchmark libraries~\cite{MarReiDua12}. Performance
is measured using the relative percentage deviation (RPD) from best-known
solutions, and statistical comparisons between algorithms are performed using
paired Wilcoxon signed-rank tests~\cite{Wilcoxon45} and Student's t-tests.

The rest of this report is organised as follows. Section~2 formally defines
the LOP. Section~3 describes the implemented algorithms and the experimental
setup. Section~4 presents and discusses the results. Section~5 concludes.
